\chapter{{\fontsize{14pt}{21pt}\selectfont\MakeUppercase{Предметная область и технологии}}}
\label{cha:analysis}

\section{Предметная область корпоративного обучения и исходное состояние}

Системы управления обучением применяются для централизованного хранения учебных материалов, назначения курсов, контроля прохождения и фиксации результатов. В корпоративной среде LMS также применяют для адаптации новых сотрудников (онбординг), обязательного обучения, внутренних регламентов, курсов по информационной безопасности и повышения квалификации сотрудников\cite{moodle_lms,open_lms}.

Для рассматриваемой предметной области характерно наличие нескольких устойчивых ролей. Системный администратор поддерживает общесистемные настройки и контролирует работу платформы. Администратор организации управляет пользователями, ролями, курсами и назначениями в пределах своей организации. Преподаватель создаёт учебный контент, тесты и анализирует результаты. Слушатель проходит назначенные курсы, изучает уроки, выполняет тесты и получает рекомендации по слабым темам.

В исходном состоянии процесс обучения часто строится с использованием нескольких несвязанных инструментов: документы и видео хранятся отдельно, списки обучающихся ведутся вручную, результаты тестирования не агрегируются, а повторение слабых тем организуется без формальной аналитики. Это приводит к ряду практических проблем:
\begin{parenlist}
  \item отсутствует единый источник данных о пользователях, курсах и результатах;
  \item повышаются трудозатраты на назначение курсов и контроль прохождения;
  \item преподаватель не получает целостной картины по ошибкам и слабым темам;
  \item слушатель проходит одинаковые по сложности тесты вне зависимости от уровня подготовки;
  \item использование мобильных устройств поддерживается ограниченно или неудобно.
\end{parenlist}

На рисунке~\ref{fig:analysis-bpmn-current} представлена схема процесса корпоративного обучения в цифровой среде. Диаграмма показывает последовательность действий от подготовки контента и назначения курса до прохождения тестирования и анализа результатов.

\begin{figure}[H]
  \centering
  \setlength{\fboxrule}{0.4pt}
  \setlength{\fboxsep}{0pt}
  \fbox{\includegraphics[width=\textwidth,height=0.82\textheight,keepaspectratio]{BPMN}}
  \caption{Схема процесса корпоративного обучения в цифровой среде}
  \label{fig:analysis-bpmn-current}
\end{figure}

\section{Постановка проблемы и назначение системы}

Проблема разработки состоит в необходимости объединить в одной системе функции администрирования, хранения контента, назначения учебных программ, мобильного доступа, тестирования и анализа результатов без усложнения пользовательских сценариев. Типовые LMS нередко обладают избыточным или трудно настраиваемым функционалом, а их мобильный пользовательский опыт не всегда соответствует сценарию коротких учебных сессий\cite{moodle_app,opigno_mobile_lms_2024}.

Назначение разрабатываемой системы заключается в поддержке полного цикла корпоративного обучения:
\begin{parenlist}
  \item создание и редактирование пользователей, ролей, курсов, уроков и тестов;
  \item назначение курсов слушателям;
  \item прохождение уроков и фиксирование прогресса;
  \item запуск адаптивного теста и подбор вопросов с учётом текущей сложности;
  \item формирование рекомендаций по слабым темам;
  \item просмотр аналитики по обучению в административном интерфейсе.
\end{parenlist}

Система должна обеспечивать работу нескольких организаций в одном развёртывании. В проекте эта задача решается не путём создания отдельной базы данных на каждого арендатора, а через логическую мультиарендную модель (несколько организаций в одном экземпляре): идентификация организации в запросе, проверка участия пользователя в организации и фильтрация данных по полю tenant\_id. Такой подход достаточен для учебного минимально жизнеспособного продукта (MVP) и соответствует текущей реализации проекта.

На рисунке~\ref{fig:analysis-usecase} показана диаграмма вариантов использования, отражающая ключевые действия ролей в системе. Диаграмма фиксирует распределение функций между системным администратором, администратором организации, преподавателем и слушателем.

\begin{figure}[H]
\centering
\includegraphics[width=\textwidth,height=0.60\textheight,keepaspectratio]{"use case"}
\caption{Диаграмма вариантов использования системы}
\label{fig:analysis-usecase}
\end{figure}

Для более формального представления сценариев в таблице~\ref{tab:analysis-usecases} приведены основные прецеденты системы, их акторы и ожидаемые результаты выполнения.

\begin{table}[H]
\centering
\caption{Ключевые прецеденты системы}
\label{tab:analysis-usecases}
{\setstretch{1}\fontsize{12pt}{12pt}\selectfont
\setlength{\tabcolsep}{5pt}
\renewcommand{\arraystretch}{1.15}
\begin{tabularx}{\textwidth}{|>{\RaggedRight\arraybackslash}p{36mm}|>{\RaggedRight\arraybackslash}p{32mm}|>{\RaggedRight\arraybackslash}X|}
\hline
{\centering\arraybackslash\textbf{Прецедент}} & {\centering\arraybackslash\textbf{Основной актор}} & {\centering\arraybackslash\textbf{Результат выполнения}} \\
\hline
Управление организациями & Системный администратор & Созданы и настроены организации, доступные для дальнейшей работы пользователей. \\
\hline
Управление пользователями и ролями & Администратор организации & В организации созданы пользователи и назначены роли в соответствии с полномочиями. \\
\hline
Управление курсами, уроками и медиа & Преподаватель & Подготовлен учебный контент, доступный для назначения слушателям. \\
\hline
Назначение курса & Администратор организации & Выбранный курс назначен слушателю и появился в его списке обучения. \\
\hline
Прохождение адаптивного теста & Слушатель & Сформирована завершённая попытка с результатом и историей ответов. \\
\hline
Получение рекомендаций & Слушатель & Пользователь получает перечень тем и материалов для повторения. \\
\hline
Просмотр аналитики & Преподаватель, администратор организации & Получены агрегированные показатели по курсам, попыткам и успеваемости. \\
\hline
\end{tabularx}
}
\end{table}

\section{Требования к системе}

Требования к системе сформированы с учётом методических рекомендаций по ВКР и подхода к анализу требований, описанного в учебнике С.~В.~Куликова. Под требованием в работе понимается проверяемое описание того, какие функции должна выполнять система и при каких условиях она должна работать, чтобы решать задачу пользователя и заказчика\cite{metodichka_vkr_ikt_2021,kulikov_testing_2024}. Поэтому требования далее разделены на функциональные и нефункциональные: первые фиксируют ожидаемое поведение программного комплекса, вторые задают ограничения на безопасность, изоляцию данных, архитектуру, интерфейсы и эксплуатационные свойства.

Бизнес-цель системы заключается в создании единого программного комплекса для корпоративного обучения, который позволяет управлять пользователями, курсами, назначениями, прохождением уроков, адаптивным тестированием и аналитикой без использования набора разрозненных инструментов. Основными ролями являются системный администратор, администратор организации, преподаватель и слушатель. Требования сформулированы так, чтобы каждое из них было атомарным, непротиворечивым, недвусмысленным и проверяемым в ходе испытаний\cite{kulikov_testing_2024}.

\subsection{Функциональные требования}

Функциональные требования описывают действия, которые система должна предоставлять пользователям и клиентским приложениям.
\begin{enumerate}
  \item FR-01 -- система должна обеспечивать регистрацию и аутентификацию пользователей.
  \item FR-02 -- система должна выдавать токены доступа и обновления для последующей работы с защищёнными функциями.
  \item FR-03 -- система должна определять активную организацию пользователя по контексту запроса.
  \item FR-04 -- система должна поддерживать роли слушателя, преподавателя, администратора организации и системного администратора.
  \item FR-05 -- администратор организации должен иметь возможность создавать пользователей и изменять состояние их доступа.
  \item FR-06 -- администратор организации должен иметь возможность назначать пользователям роли в пределах своей организации.
  \item FR-07 -- преподаватель или администратор должен иметь возможность создавать и редактировать курсы.
  \item FR-08 -- преподаватель или администратор должен иметь возможность создавать и редактировать уроки с текстовым и медийным содержимым.
  \item FR-09 -- преподаватель или администратор должен иметь возможность создавать тесты, вопросы, варианты ответов и связи вопросов с темами.
  \item FR-10 -- администратор организации должен иметь возможность назначать курс слушателю.
  \item FR-11 -- слушатель должен иметь возможность просматривать назначенные курсы, открывать уроки и фиксировать прогресс прохождения.
  \item FR-12 -- слушатель должен иметь возможность начать попытку тестирования, получать вопросы, отправлять ответы и завершать попытку.
  \item FR-13 -- после каждого ответа система должна изменять текущую сложность адаптивной попытки не более чем на один уровень.
  \item FR-14 -- система должна выбирать следующий вопрос среди ещё не заданных вопросов с учётом текущей сложности.
  \item FR-15 -- после завершения попытки система должна сохранять результат, историю ответов и выявленные слабые темы.
  \item FR-16 -- система должна формировать рекомендации по повторению материалов на основании слабых тем.
  \item FR-17 -- слушатель должен иметь возможность просматривать результат тестирования и рекомендации.
  \item FR-18 -- преподаватель и администратор организации должны иметь возможность просматривать аналитику по пользователям, курсам, тестам и назначенным обучениям.
\end{enumerate}

\subsection{Нефункциональные требования}

Нефункциональные требования фиксируют свойства качества и ограничения, которые должны учитываться при проектировании, реализации и проверке системы.
\begin{enumerate}
  \item NFR-01 -- доступ к защищённым операциям должен выполняться только после аутентификации и проверки роли пользователя,
  \item NFR-02 -- пароли пользователей должны храниться только в виде криптографических хэшей,
  \item NFR-03 -- ключевые действия пользователей должны фиксироваться в журнале аудита,
  \item NFR-04 -- данные разных организаций не должны смешиваться при выполнении пользовательских запросов,
  \item NFR-05 -- каждая сущность, зависящая от организации, должна быть связана с идентификатором организации,
  \item NFR-06 -- серверная часть должна предоставлять REST API с обменом данными в формате JSON,
  \item NFR-07 -- система должна поддерживать воспроизводимое контейнеризированное развёртывание для демонстрационного стенда,
  \item NFR-08 -- серверная часть должна быть разделена на слои и подсистемы, чтобы упростить развитие и проверку поведения,
  \item NFR-09 -- веб-интерфейс должен быть ориентирован на административные сценарии управления пользователями, курсами, тестами и аналитикой,
  \item NFR-10 -- мобильный интерфейс должен быть ориентирован на сценарии слушателя: вход, просмотр курсов, прохождение уроков, тестирование и рекомендации,
  \item NFR-11 -- пользовательский интерфейс должен поддерживать русский и английский языки,
  \item NFR-12 -- отклик основных интерактивных операций на демонстрационных данных должен быть приемлемым для работы пользователя без ощущения зависания,
  \item NFR-13 -- система должна использовать общую базу данных с логической изоляцией данных организаций,
  \item NFR-14 -- система должна хранить курсы, уроки, вопросы, попытки, результаты, слабые темы и рекомендации в структурированном виде.
\end{enumerate}

\subsection{Свойства качественных требований}

В таблице~\ref{tab:req-quality} показано, как сформулированные требования соотносятся с признаками качественных требований, описанными в учебной литературе по тестированию\cite{kulikov_testing_2024}.

\begin{table}[H]
\centering
\caption{Проверка свойств требований}
\label{tab:req-quality}
{\setstretch{1}\fontsize{12pt}{12pt}\selectfont
\setlength{\tabcolsep}{6pt}
\renewcommand{\arraystretch}{1.15}
\begin{tabularx}{\textwidth}{|>{\RaggedRight\arraybackslash}p{46mm}|>{\RaggedRight\arraybackslash}X|}
\hline
{\centering\arraybackslash\textbf{Свойство}} & {\centering\arraybackslash\textbf{Как обеспечено в работе}} \\
\hline
Завершённость & Требования охватывают основные роли, пользовательские сценарии, данные, безопасность, изоляцию организаций и клиентские приложения. \\
\hline
Атомарность & Каждое требование описывает одно проверяемое свойство или одну пользовательскую возможность. \\
\hline
Непротиворечивость & Функциональные требования не конфликтуют с ограничениями архитектуры и мультиарендной модели данных. \\
\hline
Недвусмысленность & Требования сформулированы через роли, действия пользователя и наблюдаемый результат. \\
\hline
Проверяемость & Для требований могут быть подготовлены тест-кейсы, а критический путь проверяется в программе и методике испытаний приложения~Б. \\
\hline
\end{tabularx}
}
\end{table}

\subsection{Сценарий тестирования требований по критическому пути пользователей}

Для проверки согласованности требований используется критический путь, проходящий через основные роли системы: администратор организации создаёт пользователя и назначает курс, преподаватель подготавливает курс и тест, слушатель проходит уроки и адаптивное тестирование, после чего преподаватель или администратор просматривает результаты. Такой сценарий позволяет связать требования с будущей реализацией, проектной моделью и испытаниями.

\begin{table}[H]
\centering
\caption{Трассировка требований по критическому пути}
\label{tab:req-trace}
{\setstretch{1}\fontsize{12pt}{12pt}\selectfont
\setlength{\tabcolsep}{5pt}
\renewcommand{\arraystretch}{1.15}
\begin{tabularx}{\textwidth}{|>{\RaggedRight\arraybackslash}p{32mm}|>{\RaggedRight\arraybackslash}X|>{\RaggedRight\arraybackslash}p{34mm}|}
\hline
{\centering\arraybackslash\textbf{Группа требований}} & {\centering\arraybackslash\textbf{Проверяемый пользовательский сценарий}} & {\centering\arraybackslash\textbf{Связанные требования}} \\
\hline
Пользователи и роли & Создание пользователя, вход в систему, выбор организации и проверка доступных функций & FR-01--FR-06, NFR-01--NFR-05 \\
\hline
Курсы и уроки & Создание курса, добавление урока и назначение курса слушателю & FR-07--FR-11, NFR-09, NFR-10, NFR-14 \\
\hline
Адаптивное тестирование & Запуск попытки, получение вопросов, отправка ответов, изменение сложности и завершение попытки & FR-12--FR-15, NFR-06, NFR-12 \\
\hline
Рекомендации и аналитика & Формирование рекомендаций, просмотр результата слушателем и аналитики преподавателем & FR-16--FR-18, NFR-08, NFR-14 \\
\hline
Развёртывание & Запуск демонстрационного стенда и проверка доступности серверной, веб- и мобильной частей & NFR-06--NFR-13 \\
\hline
\end{tabularx}
}
\end{table}

\section{Анализ аналогов}

Для сравнения выбраны решения, представляющие разные классы LMS: открытые платформы (\textit{Moodle}, \textit{Open LMS}), облачный вариант на базе Moodle (\textit{MoodleCloud}), платформы массовых онлайн-курсов и коммерческие корпоративные LMS\cite{moodle_lms,open_lms,moodlecloud,docebo_adaptive_2026}. В отраслевых обзорах подчёркивается значение адаптивных и интеллектуальных подходов к персонализации учебной траектории при массовом цифровом обучении\cite{educause_adaptive_2016,discovery_intelligent_adaptive}.

Сравнение выполняется не как маркетинговый обзор всех возможностей LMS, а как сопоставление решений с задачами разрабатываемого программного комплекса. Критерии выбраны из требований к системе, сформулированных выше: необходимо поддержать несколько ролей пользователей, работу слушателя с мобильного устройства, адаптивное тестирование, аналитику результатов, изоляцию данных организаций и возможность показать собственную архитектуру в рамках ВКР. Поэтому сравнение выполняется по следующим признакам:
\begin{parenlist}
  \item наличие административных ролей и разделения прав доступа --- учитывается, можно ли разграничить полномочия системного администратора, администратора организации, преподавателя и слушателя;
  \item удобство мобильного доступа для слушателя --- оценивается наличие мобильного или адаптивного интерфейса для прохождения уроков, тестов и просмотра результатов;
  \item наличие адаптивного или персонализированного тестирования --- проверяется, поддерживает ли платформа изменение сложности заданий или подбор материалов по результатам пользователя;
  \item возможности аналитики и работы со слабыми темами --- рассматривается, предоставляет ли система отчёты по успеваемости, ошибкам, темам и рекомендациям для повторения;
  \item поддержка мультиарендного сценария --- оценивается возможность вести несколько организаций в одном развёртывании с логическим разделением пользователей, курсов и результатов;
  \item пригодность для учебной разработки и доработки архитектуры --- учитывается открытость решения, возможность анализа модели данных и целесообразность использования платформы как основы выпускного проекта.
\end{parenlist}

Для формализации сравнения в таблице~\ref{tab:analysis-weights} заданы критерии и удельные веса $w_i$ такие, что $\sum_i w_i = 1$. Наибольший вес получил критерий адаптивного тестирования, поскольку именно он отличает разрабатываемую систему от типовой LMS. Мобильный доступ имеет второй по значимости вес, так как сценарий слушателя предполагает короткие учебные сессии вне рабочего места. Остальные критерии отражают обязательные условия корпоративного применения и учебной демонстрации архитектуры.

Баллы $x_{ij}$ по 10-балльной шкале выставлены на основе анализа открытых описаний, документации и функционального состава рассматриваемых решений и сведены в таблице~\ref{tab:analysis-scores}. Шкала трактуется следующим образом: 1--3 балла означают отсутствие или слабую выраженность признака, 4--6 баллов --- частичную поддержку или необходимость заметной настройки, 7--8 баллов --- наличие пригодного готового механизма, 9--10 баллов --- полное соответствие критерию в контексте целей данной ВКР. Итоговый показатель для варианта $j$ вычисляется как взвешенная сумма\cite{gost7052008}
\begin{equation}
S_j = \sum_i w_i x_{ij}.
\label{eq:analysis-score}
\end{equation}

Результаты свёртки приведены в таблице~\ref{tab:analysis-totals}. Наивысший показатель у собственной разработки, что согласуется с целью ВКР --- продемонстрировать архитектуру, модель данных и алгоритмы в открытом для анализа коде; коммерческие LMS (в таблице представлены iSpring Learn и GetCourse как ориентиры сегмента) получают высокие баллы по готовому функционалу, но низкие по критерию прозрачности архитектуры для учебного проекта. Полученные значения следует рассматривать как сравнительную оценку соответствия целям работы, а не как универсальный рейтинг LMS: крупные готовые платформы могут быть сильнее в эксплуатации и экосистеме, однако хуже подходят для демонстрации проектных решений, реализуемых в выпускной квалификационной работе.

\begin{table}[H]
\centering
\caption{Критерии сравнения LMS-аналогов и удельные веса}
\label{tab:analysis-weights}
{\setstretch{1}\fontsize{12pt}{12pt}\selectfont
\setlength{\tabcolsep}{5pt}
\renewcommand{\arraystretch}{1.15}
\begin{tabularx}{\textwidth}{|>{\RaggedRight\arraybackslash}p{32mm}|>{\RaggedRight\arraybackslash}X|>{\centering\arraybackslash}p{18mm}|}
\hline
{\centering\arraybackslash\textbf{Критерий $i$}} & {\centering\arraybackslash\textbf{Содержание}} & {\centering\arraybackslash\textbf{Вес $w_i$}} \\
\hline
$1$ & Административные роли и RBAC & 0{,}15 \\
\hline
$2$ & Мобильный доступ для слушателя & 0{,}20 \\
\hline
$3$ & Адаптивное / персонализированное тестирование & 0{,}25 \\
\hline
$4$ & Аналитика и слабые темы & 0{,}15 \\
\hline
$5$ & Мультиарендность & 0{,}15 \\
\hline
$6$ & Пригодность для анализа архитектуры в ВКР & 0{,}10 \\
\hline
\multicolumn{2}{|>{\RaggedRight\arraybackslash}X|}{Сумма весов} & 1{,}00 \\
\hline
\end{tabularx}
}
\end{table}

\begin{table}[H]
\centering
\caption{Оценка аналогов по критериям (шкала 1--10, $x_{ij}$)}
\label{tab:analysis-scores}
{\setstretch{1}\fontsize{11pt}{12pt}\selectfont
\setlength{\tabcolsep}{4pt}
\renewcommand{\arraystretch}{1.15}
\begin{tabularx}{\textwidth}{|>{\RaggedRight\arraybackslash}p{25mm}|>{\centering\arraybackslash}X|>{\centering\arraybackslash}X|>{\centering\arraybackslash}X|>{\centering\arraybackslash}X|>{\centering\arraybackslash}X|>{\centering\arraybackslash}X|}
\hline
{\centering\arraybackslash\textbf{Критерий}} & \thead{Moodle} & \thead{Open LMS} & \thead{Moodle\\Cloud} & \thead{iSpring} & \thead{Get\\Course} & \thead{Собст-\\венная} \\
\hline
$1$ & 8 & 8 & 7 & 8 & 8 & 8 \\
\hline
$2$ & 6 & 6 & 6 & 8 & 8 & 8 \\
\hline
$3$ & 5 & 5 & 5 & 7 & 7 & 9 \\
\hline
$4$ & 6 & 6 & 6 & 8 & 8 & 8 \\
\hline
$5$ & 5 & 5 & 5 & 7 & 7 & 8 \\
\hline
$6$ & 4 & 4 & 4 & 4 & 4 & 9 \\
\hline
\end{tabularx}
}
\end{table}

\begin{table}[H]
\centering
\caption{Свёртка оценок по аналогам (итоговый балл $S_j$ по формуле~\eqref{eq:analysis-score})}
\label{tab:analysis-totals}
{\setstretch{1}\fontsize{12pt}{12pt}\selectfont
\setlength{\tabcolsep}{6pt}
\renewcommand{\arraystretch}{1.15}
\begin{tabularx}{\textwidth}{|>{\RaggedRight\arraybackslash}X|>{\centering\arraybackslash}p{42mm}|}
\hline
{\centering\arraybackslash\textbf{Решение $j$}} & {\centering\arraybackslash\textbf{$S_j$ (условные единицы)}} \\
\hline
Moodle & 5{,}70 \\
\hline
Open LMS & 5{,}70 \\
\hline
MoodleCloud & 5{,}55 \\
\hline
iSpring Learn (ориентир) & 7{,}20 \\
\hline
GetCourse (ориентир) & 7{,}20 \\
\hline
Собственная разработка & 8{,}35 \\
\hline
\end{tabularx}
}
\end{table}

Поскольку административные роли поддерживаются всеми рассматриваемыми решениями, в таблице~\ref{tab:analysis-analogues} дополнительно приведено качественное сопоставление по ключевым признакам, по которым платформы различаются наиболее заметно для пользователя.

\begin{table}[H]
\centering
\caption{Сравнение аналогов по ключевым критериям}
\label{tab:analysis-analogues}
{\setstretch{1}\fontsize{10pt}{12pt}\selectfont
\setlength{\tabcolsep}{5pt}
\setlength{\extrarowheight}{3pt}
\renewcommand{\arraystretch}{1.25}
\begin{tabularx}{\textwidth}{|>{\RaggedRight\arraybackslash}m{36mm}|>{\centering\arraybackslash}m{21mm}|>{\centering\arraybackslash}m{29mm}|>{\centering\arraybackslash}m{18mm}|>{\centering\arraybackslash}X|}
\hline
\thead{Решение} &
\thead{Мобильный\\доступ} &
\thead{Адаптивное\\тестирование} &
\thead{Анали-\\тика} &
\thead{Мульти-\\арендность} \\
\hline
Moodle & средне & ограничено & да & ограничено \\
\hline
Open LMS & средне & ограничено & да & ограничено \\
\hline
MoodleCloud & средне & ограничено & да & ограничено \\
\hline
\makecell[l]{Коммерческие\\LMS} & хорошо & \makecell{частично\\или да} & хорошо & да \\
\hline
\makecell[l]{Собственная\\система} & хорошо & да & да & да \\
\hline
\end{tabularx}
}
\end{table}

По результатам сравнения можно сделать следующие выводы. Moodle и производные от него платформы обеспечивают широкий базовый функционал, однако требуют более сложной настройки и не ориентированы на демонстрацию собственной архитектуры в рамках учебной разработки. Коммерческие LMS предлагают зрелые средства мобильного доступа и аналитики, но являются закрытыми и хуже подходят как основа выпускного проекта. Разрабатываемая система уступает крупным платформам по степени зрелости, но лучше соответствует целям ВКР, поскольку позволяет явно продемонстрировать архитектуру, модель данных, адаптивное тестирование и изоляцию данных организаций-арендаторов.

Результаты сравнения показывают, что существующие решения закрывают базовые задачи LMS, однако либо обладают избыточностью и сложностью настройки, либо не ориентированы на исследование и реализацию собственной архитектуры. Для целей ВКР целесообразно разрабатывать собственный программный комплекс, в котором можно явно продемонстрировать клиент--серверную архитектуру, модель данных, алгоритм адаптивного тестирования и изоляцию данных арендаторов.

\section{Выбор технологий и платформ}

Выбор технологий определялся не только их распространённостью, но и соответствием требованиям проекта. Разрабатываемая система должна поддерживать отдельные web- и mobile-клиенты, REST API, изоляцию данных организаций-арендаторов, развитие модели предметной области и воспроизводимое локальное развёртывание. Поэтому для каждой подсистемы рассматривались несколько вариантов, после чего выбиралось решение, лучше всего соответствующее целям ВКР и текущему масштабу проекта.

\subsection{Критерии выбора технологий}

При выборе технологического стека учитывались следующие критерии:
\begin{parenlist}
  \item соответствие клиент--серверной архитектуре с отдельными web- и mobile-клиентами;
  \item удобство реализации REST API, аутентификации, разграничения доступа и контекста организации-арендатора;
  \item возможность формального описания модели данных изменений схемы БД;
  \item пригодность для минимально жизнеспособного продукта (MVP), который необходимо не только разработать, но и продемонстрировать на защите;
  \item воспроизводимость локального развёртывания без сложной инфраструктуры.
\end{parenlist}

\subsection{Сравнение технологий и результат выбора}

Сравнение основных вариантов приведено в таблице~\ref{tab:analysis-tech-choice}. В таблицу включены именно те подсистемы, которые определяют архитектурный облик проекта и влияют на удобство дальнейшей доработки.

\begin{table}[H]
\centering
\caption{Сравнение технологий и результаты выбора}
\label{tab:analysis-tech-choice}
{\setstretch{1}\fontsize{12pt}{12pt}\selectfont
\setlength{\tabcolsep}{5pt}
\renewcommand{\arraystretch}{1.15}
\begin{tabularx}{\textwidth}{|>{\RaggedRight\arraybackslash}p{32mm}|>{\RaggedRight\arraybackslash}p{45mm}|>{\RaggedRight\arraybackslash}X|}
\hline
{\centering\arraybackslash\textbf{Подсистема}} & {\centering\arraybackslash\textbf{Рассмотренные варианты}} & {\centering\arraybackslash\textbf{Выбранная технология и причина}} \\
\hline
Серверная часть & FastAPI, Django REST Framework, Flask & Выбран FastAPI, так как проект имеет формат API-first и не требует серверной генерации HTML. По сравнению с Django REST Framework решение получается легче и проще для MVP, а по сравнению с Flask требует меньше ручной сборки маршрутов, схем и документации. \\
\hline
Доступ к данным & SQLAlchemy + Alembic, Django ORM, прямой SQL & Выбраны SQLAlchemy и Alembic, поскольку они не привязаны к конкретному веб-фреймворку, позволяют явно описывать модели и контролировать миграции. Прямой SQL усложнил бы развитие, а Django ORM логичнее использовать только вместе с Django. \\
\hline
Web-клиент & React + TypeScript, Vue, Angular & Выбраны React и TypeScript. Для административной панели важны компонентный подход, маршрутизация, работа с формами и типизация данных. Angular для MVP избыточен, а Vue также подходит, но в текущем проекте более удобным оказался гибкий компонентный подход React. \\
\hline
Мобильный клиент & Flutter, React Native, native Android & Выбран Flutter, поскольку он позволяет поддерживать единый код мобильного приложения и быстро собрать интерфейс слушателя. Нативная разработка увеличила бы трудоёмкость, а React Native не давал принципиального преимущества при уже разделённых web- и mobile-сценариях. \\
\hline
Хранение данных & SQLite, PostgreSQL, MySQL & Используется комбинированный подход: SQLite удобен для локальной разработки и тестирования без отдельного сервера, а PostgreSQL лучше подходит для контейнеризированного запуска и демонстрации многопользовательской системы. \\
\hline
Развёртывание & ручной запуск сервисов, Docker Compose, Kubernetes & Выбран Docker Compose, так как он достаточен для MVP и позволяет поднимать серверную часть, web и базу данных одной конфигурацией. Ручной запуск хуже воспроизводим, а Kubernetes для учебного проекта избыточен. \\
\hline
\end{tabularx}
}
\end{table}

\subsection{Обоснование выбора серверной части}

Для серверной части рассматривались три основных подхода: использование полноценного фреймворка Django с Django REST Framework, минималистичного Flask и API-ориентированного FastAPI. В рамках данной ВКР сервер выполняет роль централизованного REST API для двух клиентов и не требует шаблонной серверной генерации страниц. По этой причине возможности Django, связанные с административной панелью, шаблонами и встроенной ORM, для проекта оказались избыточными. Flask, напротив, предоставляет слишком низкий уровень абстракции, из-за чего значительную часть инфраструктурного кода пришлось бы собирать вручную\cite{fastapi_docs,django_rest_framework}.

FastAPI оказался наиболее сбалансированным вариантом. Он хорошо подходит для явного описания маршрутов, схем запросов и ответов, зависимостей и проверки входных данных. Это соответствует фактической реализации серверной части, где сервер разделён на группы маршрутов auth, tenants, users, courses, media, tests, recommendations и analytics. Такая структура хорошо сочетается с модульным монолитом и позволяет расширять систему без переработки базовой архитектуры\cite{fastapi_docs}.

Для доступа к данным выбраны SQLAlchemy и Alembic\cite{sqlalchemy_docs}. Такой набор позволяет описывать сущности предметной области на уровне моделей, а изменения схемы фиксировать через миграции. Для ВКР это важно, поскольку система включает достаточно большую группу взаимосвязанных сущностей: организации, пользователей, роли, курсы, уроки, тесты, попытки, рекомендации и аналитические данные. Использование прямого SQL на всём протяжении проекта усложнило бы развитие и повторяемость изменений, тогда как связка SQLAlchemy + Alembic обеспечивает более управляемое развитие схемы.

\subsection{Обоснование выбора технологий клиентских приложений}

В административной части проекта рассматривались React + TypeScript, Vue и Angular. Для веб-панели требуются страницы авторизации, управления пользователями, курсами, уроками, тестами, назначениями и аналитикой. Такой интерфейс состоит из большого числа форм, таблиц и маршрутов, которые должны переиспользовать общие компоненты и работать с типизированными данными. По этой причине выбор был сделан в пользу React как гибкой компонентной библиотеки, а TypeScript был добавлен для уменьшения количества ошибок при работе с API-моделями, ролями, курсами и результатами тестирования\cite{react_docs,typescript_docs}.

Angular предоставляет более жёсткую архитектурную модель, однако для MVP это увеличивает объём шаблонного кода и порог входа. Vue также является сильным кандидатом, но в рассматриваемом проекте приоритет был отдан экосистеме React, которая хорошо сочетается с компонентной структурой административной панели, маршрутизацией и быстрым циклом сборки через Vite. Таким образом, выбор React + TypeScript оказался наиболее удобным именно для текущего состава интерфейсных сценариев.

Для мобильного клиента рассматривались Flutter, React Native и нативная Android-разработка. Так как в рамках ВКР необходимо было показать самостоятельный mobile-клиент слушателя, поддерживающий авторизацию, просмотр курсов, воспроизведение медиа, прохождение тестов и получение рекомендаций, требовалось решение с единым кодом приложения и предсказуемой сборкой интерфейса. Flutter отвечает этим требованиям и позволяет реализовать мобильный интерфейс как самостоятельную часть комплекса, не смешивая его с web-клиентом\cite{flutter_docs}.

Нативная Android-разработка на Kotlin обеспечила бы высокий уровень платформенной интеграции, но заметно увеличила бы трудоёмкость проекта. React Native позволил бы использовать близкий по идеологии стек, однако не снимал бы задачи отдельной мобильной адаптации и не давал бы существенного выигрыша по сравнению с уже выбранным web-стеком. Поэтому для данной ВКР Flutter оказался более рациональным выбором.

\subsection{Выбор СУБД и средств развёртывания}

На уровне хранения данных в проекте используется двухконтурный подход. Для локальной разработки и части тестовых сценариев применяется SQLite, так как она не требует запуска отдельного сервера, упрощает старт проекта и ускоряет проверку изменений. Для контейнеризированного развёртывания используется PostgreSQL, так как серверная СУБД лучше отражает рабочий режим многопользовательской системы, где одновременно функционируют серверная часть, web и база данных\cite{sqlite_docs,postgresql_docs}.

Такое сочетание не является противоречием. Напротив, оно позволяет использовать простой режим разработки на ранних этапах и более реалистичный режим демонстрации на этапе развёртывания. Наличие в проекте конфигураций и для SQLite, и для PostgreSQL подтверждает, что выбранный стек ориентирован как на удобство разработки, так и на последующую демонстрацию результата.

Для локального развёртывания рассматривались ручной запуск сервисов, Docker Compose и более тяжёлые оркестрационные решения. Ручной запуск серверной части, web и базы данных неудобен для повторяемой демонстрации, так как требует последовательной настройки каждого компонента. Использование Kubernetes для учебного минимально жизнеспособного продукта (MVP) не даёт соразмерной выгоды и существенно усложняет инфраструктурную часть работы. По этой причине в проекте выбран Docker Compose, который позволяет в одном файле конфигурации описать сервисы базы данных, серверной части, web и при необходимости сервис начального заполнения данными\cite{docker_compose_docs}.

Таким образом, текущий стек выбран как результат сопоставления требований проекта и трудоёмкости реализации. Он согласован со структурой каталогов исходного кода программного комплекса, обеспечивает поддержку веб- и мобильной платформ, удобен для развития предметной модели и позволяет воспроизводимо демонстрировать систему в рамках ВКР.

\section{Выводы по главе}

В первой главе рассмотрена предметная область корпоративного обучения, выявлены проблемы исходного состояния процессов и сформулировано назначение разрабатываемой системы. Определены функциональные и нефункциональные требования, приведена их проверка на основные свойства качества, а также сформирован сценарий тестирования требований по критическому пути пользователей. Выполнено сравнение аналогов и обоснован выбор собственного программного решения. На основании требований и анализа существующих платформ определён технологический стек проекта, соответствующий фактической реализации и задачам выпускной квалификационной работы.
